\section*{Agradecimientos}

A mi director, Roberto Sosa, por su guía, su confianza y la libertad que me dio para seguir las preguntas que fueron dando forma a esta tesis. A mi codirector, Alejandro D. Otero, por darme un segundo espacio donde trabajar, por sus aportes en la parte computacional y por las discusiones que enriquecieron el enfoque numérico de este trabajo.

A todo el equipo del Laboratorio de Hidrodinámica Naval y Oceánica (LabHiNO), por su trabajo en la fabricación de los modelos y en las campañas de ensayo en el canal de remolque, sin el cual la parte experimental de esta tesis no hubiera sido posible.

A Héctor R. Díaz Ojeda, por la colaboración y sus valiosos comentarios. A Stefano Gaggero y al Dipartimento di Ingegneria Navale, Elettrica, Elettronica e delle Telecomunicazioni (DITEN) de la Università di Genova, por recibirme durante mi estancia de investigación y compartir generosamente su experiencia en hélices marinas.

A la educación pública, en particular a la Universidad de Buenos Aires y a su Facultad de Ingeniería, por la formación recibida, por la beca Peruilh que sostuvo este doctorado y por la beca UBAINT Doctoral que hizo posible mi estancia en Génova. A la Iniciativa Pampa Azul y al entonces Ministerio de Ciencia, Tecnología e Innovación de la Nación, por el financiamiento otorgado.

A mis amigos y amigas del instituto, Fati, Maxi, Nico, Seba, Juancito, Mati, Fran, Cris, Oleg, Marie, e investigadores, Ceci "la jefa", Isabel, Leo, por las meriendas, la música, las discusiones y la compañía, tanto en los días en que la necesitaba como en los que no. A mis amigos y amigas de afuera de la academia, Juli, Nanu, Pao, César, [...], que esperaban verme construir un barco pero estuvieron siempre ahí.

Un agradecimiento especial a Dimas Barile, por las infinitas charlas en las que nos aconsejamos sobre áreas ajenas a las nuestras y que fueron fundamentales para destrabar numerosos problemas.

A mis padres y hermanas, por su apoyo incondicional a lo largo de todo este camino.

